\section{Evidence Preservation：让事实可以被复核}

重大事件中，团队的记忆会被高压、并行聊天和事后结果改变。证据保存不是为了“为诉讼准备故事”，而是为了让技术调查、用户保护、保险、监管回应和复盘拥有同一组可验证输入。保存动作必须尽早开始，同时避免妨碍 containment。

\subsection{Chain of custody、legal hold 与 timeline}

\term{chain of custody}（保管链）记录证据由谁、在何时、以什么方式收集、传输、访问和保存，使后来审查者能够判断完整性。Hash 可以帮助检测文件变化，但 hash 本身不证明采集方法正确。\term{legal hold}（法律保全通知）是 counsel 要求暂停常规删除并保存可能相关信息的程序；工程师不应自行宣布 privilege 或决定保全范围。

\lecturefigure{07-evidence-chain.png}{证据质量依赖 preserve、integrity、custody 与 timeline 的连续链路。}{NIST incident response guidance；本地字幕 21:00--25:40；概念重绘。}

\begin{knowledgebox}{读图：完整性来自过程，不只来自工具}
Preserve 阶段确定 logs、systems、snapshots 和 communications；integrity 记录 hash、timestamp 与 acquisition method；custody 限制并审计访问、transfer 和 purpose；timeline 把事实、假设、决策与后续修正排成时间。不能只截一张聊天截图，也不能让所有人直接修改同一证据目录。每个导出都应有 owner、来源系统、时区和保留策略。
\end{knowledgebox}

\begin{warningbox}{“清理现场”可能同时破坏证据和恢复}
删除账户、重装系统或轮换所有凭证可能是必要 containment，但若没有记录和采集，会失去 root-cause evidence。反过来，为了取证而延迟保护用户也可能扩大伤害。IC 应让 technical lead 与 forensics/counsel 明确哪些动作可立即执行、哪些需要 snapshot、哪些无法保存但必须记录理由。
\end{warningbox}

\subsection{Decision log：把事实、分析与授权分开}

高质量记录不是把每条 Slack 消息永久保存，而是维护结构化 decision log：当前技术事实、未知项、法务分析、经营选项、decision owner、时间、依据和复审条件。\term{privilege}（律师保密特权）在美国语境中通常保护为获得或提供法律意见而进行的特定保密沟通；让律师加入频道、把文档标为 privileged，并不会自动让所有业务事实消失或受保护。

\lecturefigure{09-decision-log.png}{技术事实、法务分析、高管决定与记录必须在同一链路中保持边界。}{本地字幕 21:00--25:40；NIST/SEC guidance；概念重绘。}

\begin{knowledgebox}{读图：记录“谁决定什么”，不要记录虚假的一致}
Technical facts 应包含 evidence 与 confidence；counsel analysis 写适用义务、假设和保密范围；executive decision 列出 options、risk acceptance 与 approver；record/audit 保存 owner、rationale、next review 和 superseded decision。后续事实变化时应追加修正，而不是覆盖旧条目。这样既能避免 CSO 被误认为独自作出法律判断，也能防止领导者事后声称从未收到风险信息。
\end{knowledgebox}

\begin{importantbox}{最小 decision record 模板}
\textbf{Time/owner}；\textbf{known facts and evidence links}；\textbf{unknowns}；\textbf{options}；\textbf{technical recommendation}；\textbf{legal/materiality analysis owner}；\textbf{business decision and approver}；\textbf{external communication}；\textbf{next checkpoint}。不要在事实字段写推测，也不要在 recommendation 字段伪装成已批准决定。
\end{importantbox}

\subsection{本章小结}

证据链回答“我们如何知道”，decision log 回答“谁基于哪些已知和未知作了什么决定”。两者共同保护调查质量、用户利益与组织责任，绝不是为了给个人制造免责备忘录。

